\subsection{Cấu tạo cụ thể của năm mã hóa}

Các mô tả dưới đây xét một nhóm $X=\{x_1,\ldots,x_k\}$ và chỉ mã hóa AMO.
ALO của hàng/cột luôn được thêm riêng, không phụ thuộc lựa chọn mã hóa.
Đối với Pairwise, mỗi cặp $i<j$ tạo
$(\neg x_i\lor\neg x_j)$, nên có đúng $\binom{k}{2}$ clause và không
thêm biến. Đây là lựa chọn dễ kiểm chứng nhất, nhưng trên toàn bàn cờ
số clause tăng bậc ba theo $N$ như đã tính ở trên.

Sequential Counter dùng biến phụ để biểu diễn trạng thái tiền tố đã có một
literal đúng. Những clause chuyển tiếp bảo đảm nếu một literal mới đúng thì
trạng thái tiền tố sau nó đúng; một literal mới không thể đúng khi tiền tố
trước đó đã có literal đúng. Project gọi
\texttt{CardEnc.atmost(..., encoding=EncType.seqcounter)} thay vì tự viết
chuỗi clause. Cách triển khai này có số biến phụ và clause tuyến tính theo
$k$; số lượng chính xác phụ thuộc quy tắc biên của thư viện.

Binary, tức \texttt{EncType.bitwise} trong PySAT, gán cho vị trí $i$ một
mã nhị phân có $\lceil\log_2 k\rceil$ bit phụ. Nếu $x_i$ đúng, các bit
phải khớp mã của $i$. Hai literal khác nhau không thể đồng thời đúng vì
chúng đòi ít nhất một bit nhận hai giá trị trái nhau. Kết quả là số bit
tăng theo $\log k$, trong khi số clause tăng theo $k\log k$. Mã nhị phân
không đồng nghĩa với ít clause nhất; nó đánh đổi số biến phụ lấy các
liên hệ giữa literal gốc và bit mã.

Commander do project tự cài đặt. Đặt $g=\lceil\sqrt{k}\rceil$, chia nhóm
thành $m=\lceil k/g\rceil$ nhóm có kích thước $g_1,\ldots,g_m$.
Mỗi nhóm có một biến chỉ huy $c_j$. Các clause
$(\neg x\lor c_j)$ cho từng thành viên và
$(\neg c_j\lor\bigvee_{x\in G_j}x)$ buộc $c_j$ đúng đúng khi nhóm
có thành viên đúng. Pairwise AMO được dùng cả bên trong nhóm và trên
các chỉ huy. Với nhánh áp dụng khi $k>3$, số biến phụ là $m$ và số
clause theo đúng implementation là
\begin{equation}
 C_{\mathrm{cmd}}(k)
 =\sum_{j=1}^{m}\binom{g_j}{2}+\binom{m}{2}+k+m.
\end{equation}
Ở các nhóm kích thước không quá ba, code chuyển thẳng sang Pairwise.

Product cũng do project tự cài đặt. Đặt
$p=\lceil\sqrt{k}\rceil$ và $q=\lceil k/p\rceil$, rồi gắn mỗi literal
với một ô khác nhau trong lưới $p\times q$. Nếu literal đúng, hai chỉ
báo hàng và cột của ô đó phải đúng. Pairwise AMO trên các chỉ báo hàng
và cột khiến không thể có hai literal gốc đúng. Với $k>4$, có $p+q$
biến phụ và đúng
\begin{equation}
 C_{\mathrm{prod}}(k)
 =2k+\binom{p}{2}+\binom{q}{2}
\end{equation}
clause. Với $k\leq4$, implementation dùng Pairwise trực tiếp.

Các công thức cho thấy bậc tiệm cận chưa quyết định lựa chọn ở nhóm nhỏ.
Ví dụ $k=8$: Pairwise tạo 28 clause; Commander chia nhóm kích thước
$3,3,2$ và tạo 21 clause cùng ba biến phụ; Product dùng lưới $3\times3$
và tạo 22 clause cùng sáu biến phụ. Đây là số cho một nhóm AMO, chưa
tính clause ALO hay nhiều nhóm chồng lấn của toàn bàn cờ.

\subsection{Chi phí ở quy mô toàn bàn cờ}

Có $O(N)$ nhóm AMO độ dài tối đa $N$ và tổng độ dài các nhóm là
$O(N^2)$. Do đó Pairwise cho $O(N^3)$ clause; Sequential và Product
cho $O(N^2)$; Binary cho $O(N^2\log N)$; Commander cài đặt ở đây
cho $O(N^{5/2})$. Đây là bậc trên theo cách dựng công thức, không phải
độ phức tạp thời gian tìm kiếm của Glucose3. Ngay cả khi hai encoding
có cùng bậc clause, số literal mỗi clause và cấu trúc lan truyền có thể
khác nhau. Vì vậy phần thực nghiệm đo cả số biến, số clause, thời gian
và run hoàn thành thay vì chọn encoding bằng phân tích tiệm cận đơn thuần.
